\chapter{Prototype Integration and Experimental Methods}\label{ch:prototype-integration-and-experimental-methods}

This chapter separates completed subsystem bring-up from the
experimental work still required to validate an integrated, freely
balancing robot. Repository evidence and project records support
operation of the four actuators, IMU acquisition and estimation, CRSF
receiver decoding, display output, CAN\slash{}TWAI communication, and
supervisor infrastructure. They do not yet support a claim of
closed-loop balancing or omnidirectional motion. Consequently, Sections
10.1 and 10.2 document the present integration baseline, while the
remaining sections specify proposed procedures in future tense.

\begin{longtable}[]{@{}
  >{\raggedright\arraybackslash}p{(\columnwidth - 0\tabcolsep) * \real{1.0000}}@{}}
\toprule\noalign{}
\begin{minipage}[b]{\linewidth}\raggedright
\textless INSERT FIGURE 10.1: Photograph the assembled robot on a
neutral background. Label the four collinear Mecanum wheels, actuator
modules, frame, battery, controller, IMU, CAN harness, receiver,
display, center-of-mass region, balancing axis, and body axes. Include a
metric scale and image date\textgreater{}\\
What the figure should show: Photograph the assembled robot on a neutral
background. Label the four collinear Mecanum wheels, actuator modules,
frame, battery, controller, IMU, CAN harness, receiver, display,
center-of-mass region, balancing axis, and body axes. Include a metric
scale and image date\\
Likely source: dated photograph of the final physical prototype\strut
\end{minipage} \\
\midrule\noalign{}
\endhead
\bottomrule\noalign{}
\endlastfoot
\end{longtable}

\begin{center}\singlespacing
\phantomsection\addcontentsline{lof}{figure}{\protect\numberline{10.1}Integrated prototype prepared for controlled bench testing.}
\textbf{Figure 10.1. Integrated prototype prepared for controlled bench testing.}
\end{center}

\section{Integration Baseline}\label{integration-baseline}

The confirmed development platform uses an ESP32-class controller and
FreeRTOS-based firmware. The four selected AK40-10 actuators can be
commanded during bring-up, and the firmware contains four actuator
slots; however, the inspected configuration explicitly identified only
one active node address. The ISM330DHCX driver reads a contiguous sample
containing accelerometer and gyroscope data, applies a mounting
transform, calibrates gyro bias, and supplies a complementary attitude
estimator. Receiver and display paths have also been exercised. These
achievements establish interface readiness but not timing determinism,
calibrated dynamics, or closed-loop stability.

\clearpage
\begin{landscape}
\par\medskip\noindent\singlespacing
\phantomsection\addcontentsline{lot}{table}{\protect\numberline{10.1}Integration and experimental test matrix}
\textbf{Table 10.1. Integration and experimental test matrix}\par

\begin{longtable}[]{@{}
  >{\raggedright\arraybackslash}p{(\columnwidth - 8\tabcolsep) * \real{0.0846}}
  >{\raggedright\arraybackslash}p{(\columnwidth - 8\tabcolsep) * \real{0.2692}}
  >{\raggedright\arraybackslash}p{(\columnwidth - 8\tabcolsep) * \real{0.2077}}
  >{\raggedright\arraybackslash}p{(\columnwidth - 8\tabcolsep) * \real{0.2615}}
  >{\raggedright\arraybackslash}p{(\columnwidth - 8\tabcolsep) * \real{0.1769}}@{}}
\toprule\noalign{}
\begin{minipage}[b]{\linewidth}\raggedright
\textbf{ID}
\end{minipage} & \begin{minipage}[b]{\linewidth}\raggedright
\textbf{Test}
\end{minipage} & \begin{minipage}[b]{\linewidth}\raggedright
\textbf{Present status}
\end{minipage} & \begin{minipage}[b]{\linewidth}\raggedright
\textbf{Acceptance evidence}
\end{minipage} & \begin{minipage}[b]{\linewidth}\raggedright
\textbf{Disposition}
\end{minipage} \\
\midrule\noalign{}
\endhead
\bottomrule\noalign{}
\endlastfoot
I-01 & Mechanical assembly and wheel-order survey & Partially documented
& Dimensioned drawing, photographs, torque record & Complete before
modeling freeze \\
I-02 & Power-off continuity and insulation & No record located & Signed
inspection sheet & Pending \\
I-03 & Current-limited regulator bring-up & No final-PCB record located
& Rail voltage\slash{}ripple plots & Pending \\
I-04 & CAN physical layer and four-node enumeration & Single-node
evidence; four slots in code & Oscilloscope trace and node table &
Pending \\
I-05 & IMU identity, axes, bias, and time stamps & Driver\slash{}estimator
implemented & Calibration record and sign tests & Repeat on final
mount \\
I-06 & Receiver failsafe and display & Subsystem operation reported &
Recorded timeout and UI test & Document \\
I-07 & Restrained balance loop & Not confirmed & Synchronized
closed-loop log & Pending \\
I-08 & Free balance and disturbance rejection & Not confirmed & Repeated
trials and metrics & Pending \\
I-09 & Longitudinal\slash{}lateral\slash{}yaw motion & Not confirmed & Reference and
measured trajectories & Pending \\
I-10 & Power, thermal, endurance, and fault tests & Not confirmed &
Logged voltage\slash{}current\slash{}temperature\slash{}errors & Pending \\
\end{longtable}
\end{landscape}
\clearpage

\emph{Note.} A completed subsystem test does not satisfy a vehicle-level
requirement unless its final hardware, configuration, and interfaces are
represented.

\section{Configuration Control and Pre-Test Record}\label{configuration-control-and-pre-test-record}

Every trial should be associated with a configuration record containing
the repository commit, firmware build flags, parameter-file checksum,
controller board revision, PCB revision, actuator firmware and operating
mode, wheel order and handedness, battery identification and state of
charge, payload, tire condition, and raw-data filename. Configuration
control is necessary because a sign change, node reassignment, filter
coefficient, or current limit can materially change the outcome.

\begin{longtable}[]{@{}
  >{\raggedright\arraybackslash}p{(\columnwidth - 0\tabcolsep) * \real{1.0000}}@{}}
\toprule\noalign{}
\begin{minipage}[b]{\linewidth}\raggedright
\textless VERIFY FROM HARDWARE: Confirm the installed motor model and
quantity; firmware version; actuator command mode; four CAN node IDs;
CAN bitrate; current, voltage, speed, temperature, and ERPM limits;
encoder zero convention; wheel-direction signs; radio failsafe values;
controller build identifier; and final PCB revision before
testing.\textgreater{}
\end{minipage} \\
\midrule\noalign{}
\endhead
\bottomrule\noalign{}
\endlastfoot
\end{longtable}

\clearpage
\begin{landscape}
\par\medskip\noindent\singlespacing
\phantomsection\addcontentsline{lot}{table}{\protect\numberline{10.2}Test instrumentation and measurement channels}
\textbf{Table 10.2. Test instrumentation and measurement channels}\par

\begin{longtable}[]{@{}
  >{\raggedright\arraybackslash}p{(\columnwidth - 8\tabcolsep) * \real{0.1825}}
  >{\raggedright\arraybackslash}p{(\columnwidth - 8\tabcolsep) * \real{0.2482}}
  >{\raggedright\arraybackslash}p{(\columnwidth - 8\tabcolsep) * \real{0.1606}}
  >{\raggedright\arraybackslash}p{(\columnwidth - 8\tabcolsep) * \real{0.1168}}
  >{\raggedright\arraybackslash}p{(\columnwidth - 8\tabcolsep) * \real{0.2920}}@{}}
\toprule\noalign{}
\begin{minipage}[b]{\linewidth}\raggedright
\textbf{Quantity}
\end{minipage} & \begin{minipage}[b]{\linewidth}\raggedright
\textbf{Instrument\slash{}channel}
\end{minipage} & \begin{minipage}[b]{\linewidth}\raggedright
\textbf{Range\slash{}resolution}
\end{minipage} & \begin{minipage}[b]{\linewidth}\raggedright
\textbf{Rate}
\end{minipage} & \begin{minipage}[b]{\linewidth}\raggedright
\textbf{Calibration\slash{}traceability}
\end{minipage} \\
\midrule\noalign{}
\endhead
\bottomrule\noalign{}
\endlastfoot
Pitch and pitch rate & Onboard IMU plus independent reference &
\textless INSERT\textgreater{} & \textless INSERT Hz\textgreater{} &
Six-position accelerometer and stationary gyro tests \\
Wheel position\slash{}speed & Four actuator feedback channels &
\textless INSERT\textgreater{} & \textless INSERT Hz\textgreater{} &
Direction and counts\slash{}revolution verification \\
Motor current & Controller feedback plus external sensor &
\textless INSERT A\textgreater{} & \textless INSERT Hz\textgreater{} &
Compare at two or more known loads \\
Battery voltage & Onboard ADC\slash{}BMS plus DMM & \textless INSERT
V\textgreater{} & \textless INSERT Hz\textgreater{} & Two-point
comparison \\
Body pose & Motion capture, camera, or surveyed floor markers &
\textless INSERT\textgreater{} & \textless INSERT Hz\textgreater{} &
Scale and coordinate-frame calibration \\
Temperatures & Actuator\slash{}controller\slash{}battery sensors & \textless INSERT
\ensuremath{^{\circ}}C\textgreater{} & \textless INSERT Hz\textgreater{} & Ambient and
reference comparison \\
Loop timing\slash{}CAN & Firmware time stamps and bus analyzer &
\textless INSERT \ensuremath{\mu}s\textgreater{} & Per event & Common time base or
measured offset \\
\end{longtable}
\end{landscape}
\clearpage

\begin{longtable}[]{@{}
  >{\raggedright\arraybackslash}p{(\columnwidth - 0\tabcolsep) * \real{1.0000}}@{}}
\toprule\noalign{}
\begin{minipage}[b]{\linewidth}\raggedright
\textless INSERT FIGURE 10.2: Show a bench power supply or protected
battery connection, current limit, fuses, emergency disconnect,
controller, CAN analyzer\slash{}oscilloscope, four actuators, termination
points, and measurement references. Label expected voltage domains and
the safe operator position\textgreater{}\\
What the figure should show: Show a bench power supply or protected
battery connection, current limit, fuses, emergency disconnect,
controller, CAN analyzer\slash{}oscilloscope, four actuators, termination
points, and measurement references. Label expected voltage domains and
the safe operator position\\
Likely source: laboratory photograph and annotated electrical
diagram\strut
\end{minipage} \\
\midrule\noalign{}
\endhead
\bottomrule\noalign{}
\endlastfoot
\end{longtable}

\begin{center}\singlespacing
\phantomsection\addcontentsline{lof}{figure}{\protect\numberline{10.2}Current-limited electrical bring-up and CAN validation setup.}
\textbf{Figure 10.2. Current-limited electrical bring-up and CAN validation setup.}
\end{center}

\section{Power-On and Communication Validation}\label{power-on-and-communication-validation}

With the wheels unloaded and propulsion disabled, the proposed procedure
begins with resistance and polarity checks. Logic rails are energized
from a current-limited source before actuator power is connected. Rail
voltage, ripple, regulator temperature, and idle current are recorded.
CANH and CANL resistance is measured with power removed; approximately
60 \ensuremath{\Omega} is expected only when two 120 \ensuremath{\Omega} end terminations are present in
parallel \cite{ti2026tida01238}. With power applied, recessive and dominant levels,
edge integrity, error counters, and traffic load are observed. Each
motor is then enabled individually at a low command while its physical
direction and feedback sign are recorded.

Abort criteria include reversed polarity, unexpected current, regulator
dropout, repeated CAN error-passive or bus-off transitions, actuator
faults, uncontrolled motion, excessive temperature, stale feedback, or
disagreement between commanded and observed direction. The emergency
disconnect must be reachable without crossing the
robot\textquotesingle s fall envelope.

\section{Sensor Calibration and Coordinate-Frame Verification}\label{sensor-calibration-and-coordinate-frame-verification}

The final IMU installation will be calibrated after the mounting
hardware is fixed. Accelerometer offset and scale are estimated using
six static orientations; gyroscope bias and noise are estimated from a
stationary record after thermal stabilization. The body-to-sensor
rotation is then verified by applying positive roll, pitch, and yaw
rotations one at a time. A forward chassis tilt must produce the pitch
sign assumed by the controller, and a positive wheel command must
produce the modeled body-force sign. Independent observation is required
because an internally consistent but inverted sensor-and-actuator
convention can pass software unit tests while destabilizing the physical
robot.

\clearpage
\begin{landscape}
\par\medskip\noindent\singlespacing
\phantomsection\addcontentsline{lot}{table}{\protect\numberline{10.3}Calibration record}
\textbf{Table 10.3. Calibration record}\par

\begin{longtable}[]{@{}
  >{\raggedright\arraybackslash}p{(\columnwidth - 8\tabcolsep) * \real{0.2000}}
  >{\raggedright\arraybackslash}p{(\columnwidth - 8\tabcolsep) * \real{0.2231}}
  >{\raggedright\arraybackslash}p{(\columnwidth - 8\tabcolsep) * \real{0.1692}}
  >{\raggedright\arraybackslash}p{(\columnwidth - 8\tabcolsep) * \real{0.2231}}
  >{\raggedright\arraybackslash}p{(\columnwidth - 8\tabcolsep) * \real{0.1846}}@{}}
\toprule\noalign{}
\begin{minipage}[b]{\linewidth}\raggedright
\textbf{Item}
\end{minipage} & \begin{minipage}[b]{\linewidth}\raggedright
\textbf{Method}
\end{minipage} & \begin{minipage}[b]{\linewidth}\raggedright
\textbf{Raw record}
\end{minipage} & \begin{minipage}[b]{\linewidth}\raggedright
\textbf{Result}
\end{minipage} & \begin{minipage}[b]{\linewidth}\raggedright
\textbf{Date\slash{}configuration}
\end{minipage} \\
\midrule\noalign{}
\endhead
\bottomrule\noalign{}
\endlastfoot
Accelerometer bias\slash{}scale & Six static orientations & \textless INSERT
FILE\textgreater{} & \textless INSERT vector\slash{}matrix\textgreater{} &
\textless INSERT\textgreater{} \\
Gyro bias\slash{}noise & Stationary thermal record & \textless INSERT
FILE\textgreater{} & \textless INSERT rad\slash{}s and noise\textgreater{} &
\textless INSERT\textgreater{} \\
IMU mounting rotation & Single-axis sign tests & \textless INSERT
FILE\textgreater{} & \textless INSERT transform\textgreater{} &
\textless INSERT\textgreater{} \\
Wheel direction and feedback & Lifted low-speed command &
\textless INSERT FILE\textgreater{} & \textless INSERT four
signs\textgreater{} & \textless INSERT\textgreater{} \\
Effective radius & Loaded rollout & \textless INSERT FILE\textgreater{}
& \textless INSERT m \ensuremath{\pm} uncertainty\textgreater{} &
\textless INSERT\textgreater{} \\
Current and voltage channels & Reference instrument comparison &
\textless INSERT FILE\textgreater{} & \textless INSERT
calibration\textgreater{} & \textless INSERT\textgreater{} \\
\end{longtable}
\end{landscape}
\clearpage

\section{Restrained Controller Commissioning}\label{restrained-controller-commissioning}

Initial balance trials will use a low-friction restraint or pivot that
prevents a fall from damaging equipment without carrying normal balance
torque. The system will start in DISARMED, verify fresh IMU and actuator
feedback, confirm a valid receiver state, and require an explicit
neutral arming command. The first closed-loop test will use conservative
torque and slew limits with integral action disabled. Feedback polarity
will be checked using a small manual displacement while the wheels are
clear of the ground, followed by low-authority ground contact.
Controller bandwidth will be increased only after delay, noise,
saturation, and structural compliance have been measured.

\begin{longtable}[]{@{}
  >{\raggedright\arraybackslash}p{(\columnwidth - 0\tabcolsep) * \real{1.0000}}@{}}
\toprule\noalign{}
\begin{minipage}[b]{\linewidth}\raggedright
\textless INSERT FIGURE 10.3: Depict the restraint geometry, allowed
pitch motion, load-free directions, emergency stop, operator position,
protected fall envelope, and any tether forces. The restraint must not
create an unmodeled restoring moment over the tuning
range\textgreater{}\\
What the figure should show: Depict the restraint geometry, allowed
pitch motion, load-free directions, emergency stop, operator position,
protected fall envelope, and any tether forces. The restraint must not
create an unmodeled restoring moment over the tuning range\\
Likely source: CAD drawing and laboratory photograph of the
commissioning fixture\strut
\end{minipage} \\
\midrule\noalign{}
\endhead
\bottomrule\noalign{}
\endlastfoot
\end{longtable}

\begin{center}\singlespacing
\phantomsection\addcontentsline{lof}{figure}{\protect\numberline{10.3}Mechanical restraint used for initial balance-controller tuning.}
\textbf{Figure 10.3. Mechanical restraint used for initial balance-controller tuning.}
\end{center}

\section{Free-Balance and Disturbance-Rejection Tests}\label{free-balance-and-disturbance-rejection-tests}

After restrained stability is demonstrated, free-balance trials will be
performed within a clear padded area. A test begins from a documented
pitch angle and neutral planar command. At least five successful repeats
and all unsuccessful attempts will be retained. Pitch settling time,
percent overshoot where meaningful, steady-state bias, RMS error over a
stated interval, peak command, saturation duration, and feedback age are
computed using fixed definitions. Disturbance rejection will use a
repeatable impulse mechanism or a controlled initial displacement; an
undocumented hand push is acceptable only as a qualitative demonstration
and not as the principal quantitative result.

\begin{longtable}[]{@{}
  >{\raggedright\arraybackslash}p{(\columnwidth - 0\tabcolsep) * \real{1.0000}}@{}}
\toprule\noalign{}
\begin{minipage}[b]{\linewidth}\raggedright
\textless INSERT DATA: Measured pitch-angle response during recovery
from an approximately [ANGLE]-degree initial displacement. Required
signals: pitch angle (deg), pitch rate (deg\slash{}s), four motor commands and
feedback values, saturation flags, operating state, IMU age,
CAN-feedback age, battery voltage (V), and time (s). Provide at least
five trials and raw machine-readable files.\textgreater{}
\end{minipage} \\
\midrule\noalign{}
\endhead
\bottomrule\noalign{}
\endlastfoot
\end{longtable}

\section{Longitudinal, Lateral, and Yaw Tests}\label{longitudinal-lateral-and-yaw-tests}

Planar tests will be introduced sequentially while the balance loop
remains active. Step and trapezoidal references will be used for
longitudinal velocity, lateral velocity, and yaw rate, followed by
combined commands. Reference magnitudes will be bounded by measured
torque, speed, traction, and balance reserves. External pose measurement
is preferred because wheel integration cannot distinguish true lateral
motion from roller slip. Each trial will report reference and measured
motion, pitch response, all wheel commands, current, saturation, and the
allocation residual.

The collinear geometry warrants an additional identifiability check:
pure lateral and pure yaw commands must generate measurably distinct
wheel-speed patterns and body responses. If H is poorly conditioned or
the assembled roller pattern is rank deficient, the affected degree of
freedom will be reported as constrained rather than hidden through
controller tuning.

\section{Power, Thermal, Endurance, and Fault Tests}\label{power-thermal-endurance-and-fault-tests}

A staged endurance test will log battery voltage, branch current,
regulator voltage and temperature, actuator temperatures, CAN errors,
loop overruns, and fault state while the robot executes a repeatable
bounded maneuver. Test duration and ambient conditions must be stated.
Separate injected-fault tests will remove receiver updates, delay motor
feedback, disconnect one actuator while safe, exceed the permitted tilt,
and approach undervoltage thresholds. The expected response is a
deterministic transition to a documented safe state, not merely
cessation of logging. Battery testing will follow the battery
manufacturer\textquotesingle s handling and discharge limits.

\begin{longtable}[]{@{}
  >{\raggedright\arraybackslash}p{(\columnwidth - 0\tabcolsep) * \real{1.0000}}@{}}
\toprule\noalign{}
\begin{minipage}[b]{\linewidth}\raggedright
\textless INSERT FIGURE 10.4: Show the shared time base among firmware
logs, CAN analyzer, external pose or video, current\slash{}voltage
instrumentation, and thermal sensors. Identify trigger events, measured
synchronization error, sample rates, and file naming\textgreater{}\\
What the figure should show: Show the shared time base among firmware
logs, CAN analyzer, external pose or video, current\slash{}voltage
instrumentation, and thermal sensors. Identify trigger events, measured
synchronization error, sample rates, and file naming\\
Likely source: data-acquisition schematic and photograph\strut
\end{minipage} \\
\midrule\noalign{}
\endhead
\bottomrule\noalign{}
\endlastfoot
\end{longtable}

\begin{center}\singlespacing
\phantomsection\addcontentsline{lof}{figure}{\protect\numberline{10.4}Synchronized data-acquisition and reference-measurement arrangement.}
\textbf{Figure 10.4. Synchronized data-acquisition and reference-measurement arrangement.}
\end{center}

\section{Repeatability, Uncertainty, and Data Reduction}\label{repeatability-uncertainty-and-data-reduction}

Raw data will remain immutable and will be processed by
version-controlled scripts. Time stamps will be converted to seconds on
a common origin, duplicate or missing samples will be reported, and
filtering will be applied only once with documented order, cutoff, and
phase characteristics. Uncertainty will combine sensor calibration, time
alignment, repeatability, and reference-system resolution where
applicable. Means, standard deviations, confidence intervals, and trial
counts will accompany aggregate metrics. Failed trials are part of the
evidence and will not be discarded without a stated exclusion rule.

\begin{longtable}[]{@{}
  >{\raggedright\arraybackslash}p{(\columnwidth - 0\tabcolsep) * \real{1.0000}}@{}}
\toprule\noalign{}
\begin{minipage}[b]{\linewidth}\raggedright
\textless INSERT PROCEDURE: Provide the final laboratory risk
assessment, approved battery-charging and storage procedure,
emergency-stop design, test-area layout, required personal protective
equipment, maximum authorized command\slash{}current\slash{}tilt limits, abort
criteria, and names of trained observers.\textgreater{}
\end{minipage} \\
\midrule\noalign{}
\endhead
\bottomrule\noalign{}
\endlastfoot
\end{longtable}
